\documentclass{article}
\usepackage{amsmath,hyperref}
\begin{document}
\section*{Public phase history can trade time for writer evidence and authority}
\paragraph{Matched scope.}
Use the same two swapped residual gates, actual Open-RMF phase
and writer excerpts, serial velocity-following plant and strict
aperture margin as
\href{https://github.com/eternal-november/eos-research-programmes/pull/828}{the expiring-hold fixture}.
Four updates are required for passage. The late CLOSE, if any,
first becomes eligible at tick 1, targets exactly one fixed gate,
and cannot be applied while its authored hold lasts. No real
ROS delay, stock hold API, body dynamics, or field result is assumed.
The old fast comparator still exits at 5 using a writer read and
expiry at 2. This study varies the public observation cut and
delivery separately up to a deadline of 12, holding the platform,
clearance, primitive histories and terminal service predicate fixed.

\paragraph{No hold and no privileged read.}
If the high-residual gate is the first to leave OPEN at tick 2,
it received CLOSE, and the other gate is the usable one.
If the high-residual gate first leaves at tick 3, use it after
it recovers. The low-residual gate does not leave OPEN by tick 3
in either case. Retaining each gate's first non-OPEN boundary
through tick 3 suffices, as does retaining the full pair of
ordered sequences through that boundary. The data may arrive as
late as tick 7. At 7 the chosen gate can be traversed by 11.
When CLOSE targets the low gate, no crossing of either gate
can start before 7, even knowing the history. Both observations
through tick 2 can match in the opposite low-CLOSE histories,
which force opposite gate decisions. A tick-3 or tick-6 snapshot
does not replace the ordered prefix for arrival at 7. A fresh
tick-7 snapshot works at 7, but if it arrives at tick 8, the
exit shifts to 12. Delivery at 8 of even complete prior evidence
cannot beat the four-update crossing bound with deadline 11.

\paragraph{A protective change can remove usable evidence.}
A mandatory hold through tick 1 is stricter than the no-hold
rule on which the old tick-3 history policy relies. With expiry
at 2 and the same report deadline and release cut, it delays
the high-gate CLOSE enough to mask the tick-2 signature. For a
fixed high gate, opposite CLOSE destinations now produce the
same early phase history, but at entry 7 one admits only the
low gate, and the other only the high gate. The early
public-history policy no longer exists. This is not a claim
that permission to install an optional hold reduces an
optimizer's choices; the policy may decline the hold. It
shows why an intervention changing the evidence process is
not just a restriction of the old action menu.

\paragraph{Other matched capability points.}
A simultaneous phase pair taken at 3 suffices to exit at 7 with
expiry at 4. Delivering that pair at 4 requires expiry at 5 in
the checked positive arm and exits at 8. A fixed gate can be
used from 7 to 11 without an observation when the hold lasts
until expiry at 8. The same fixed-gate interface at zero hold
does not suffice by 11. These points are conditional capability
tradeoffs, not measured installation or communication costs.

\paragraph{Verification.}
The native extraction actually publishes and evaluates each
phase, writes actuator velocity and continues an unrelated third
writer. The plant integrates each velocity one update at a time;
the test checks all five boundaries of the crossing. Every native
aperture trace matches an independent integer recurrence for the
stated plant. The fixed policy compiler and exhaustive same-view
direct planner agree on all 16 arms; the compiled plans are rerun.
Six Lean checks establish the finite delay, scope and
indistinguishability witnesses, not a general simulation soundness
theorem. No independent mission or architectural gain is asserted.
\end{document}
